\paragraph{LCP2: correcting for the variance of the leak estimate (plan frozen before the run).} Theorem~\ref{thm:lcp}(b) identifies a null variance term of order $1/n_d$ that LCP's scaling omits: the variance of the per-donor leak estimate and its covariance with $r_{dj}r_{dk}$. LCP2 estimates it per donor by a delete-one-nucleus jackknife, $\hat\delta_d=\widehat{\mathrm{Var}}(r_{dj}r_{dk}-\tilde\ell_d)-\widehat{\mathrm{Var}}(r_{dj}r_{dk})$, and adds $\max(\sum_d\hat\delta_d,0)$ to the denominator of the scaling factor; the permutation reference is unchanged. On new seeds and the held-out gene panel ($R=100$ per design), LCP2's FDR was 0.052\,(0.018) at 16 donors with ${\sim}40$ cells (LCP 0.110) and 0.025\,(0.003) with the Micro-PVM\_1 cell numbers (LCP 0.077), and stayed below 0.05 in the other four designs (0.033--0.042). In the real-cell null it made no calls and its per-pair type-I error was 0.046 and 0.050, but 0.068 and 0.085 for the 1\% most within-correlated pairs, so LCP2, like LCP, fails the frozen real-cell criterion (top-1\% type-I at most $0.05+3$\,MCSE $\approx0.060$). The price is power: 0.106 versus 0.122 at 24 donors (0.87 of LCP's), 0.482 versus 0.501 at 40, and 0.298 versus 0.399 with Micro-PVM\_1 cell numbers. Our frozen rule required at least 0.9 of LCP's power at 24 and 40 donors, so LCP2 does not replace LCP; as a post hoc recommendation, we suggest LCP2 for designs like the two where LCP failed (16 donors with ${\sim}40$ cells; ${\sim}15$ cells per donor), accepting its lower power.
